% !TEX root = ../ComputationalOTFnT.tex
\newcommand{\myFigRKHS}[1]{\includegraphics[width=.195\linewidth]{rkhs/#1}}
\chapter{統計的ダイバージェンス}
\label{c-statistical}

本章では、Wasserstein距離の統計的性質を考察する。より具体的には、データサイエンスで日常的に用いられる他の主要な距離やダイバージェンスとの比較を行う。二つの確率分布間の距離を、それらの分布からのサンプルのみを用いてどのように近似できるかを定量化する。これらの話題を導入するために、\S\ref{sec-phi-div}と\S\ref{sec-dual-norms}では、それぞれ確率分布間のダイバージェンスと積分確率計量について概説する。ダイバージェンス$D$は典型的には$D(\al,\be) \geq 0$かつ$D(\al,\be)=0$であることと$\al=\be$であることが同値であるが、対称性や三角不等式を満たす必要はない。測度に対する積分確率計量は、所定のテスト関数族を用いて定義される双対ノルムである。これらの量はWasserstein距離に対する健全な代替であり、\S\ref{c-variational}で扱う推論問題に取り組むための損失関数として日常的に用いられている。まず\S\ref{sec-non-embeddability}において、最適輸送距離がHilbert的でないこと、すなわち確率測度の適切な特徴表現上のHilbert計量を用いて効率的に近似することができないことを示す。\S\ref{sec-empirical-wass}では、$\al$と$\be$からそれぞれ抽出された離散サンプル$(x_i)_i$と$(y_j)_j$から$D(\al,\be)$をどのように近似するかを示す。この問題に対する十分な統計的理解は、機械学習においてWasserstein距離を用いる際に不可欠である。なお、本節では主に\emph{連続}集合上に台を持つ分布間の最適輸送の統計的近似を扱う。台空間が有限である場合の同じ問題は、~\citet{sommerfeld2018inference}の研究に続き文献で注目を集めており、~\citet{bigot2017central}によってエントロピー正則化量に拡張されている。


%%%%%%%%%%%%%%%%%%%%%%%%%%%%%%%%%%%%%%%%%%%%%%%%%%%%%%%%%%%%%%%%%%%%%%%%%%%%%%%%%%%%%%%%%%%%%%%%%%%%
%%%%%%%%%%%%%%%%%%%%%%%%%%%%%%%%%%%%%%%%%%%%%%%%%%%%%%%%%%%%%%%%%%%%%%%%%%%%%%%%%%%%%%%%%%%%%%%%%%%%
%%%%%%%%%%%%%%%%%%%%%%%%%%%%%%%%%%%%%%%%%%%%%%%%%%%%%%%%%%%%%%%%%%%%%%%%%%%%%%%%%%%%%%%%%%%%%%%%%%%%
\section{$\phi$-ダイバージェンス}
\label{sec-phi-div}

次節で$\Wass_1$と類似した構成を持つ「弱い」ノルムを詳述する前に、いわゆるダイバージェンスの一般的な構成について説明する。これは確率分布を推定する際の損失関数として利用できる。このようなダイバージェンスは、質量の移動の概念を導入することなく、二つの入力測度をそれらの質量の\emph{各点}比較によって比較する。ダイバージェンスは、二つの測度間の各点の比率を見ることで、それらがどれほど近いかの感覚を与える汎関数である。これらは良好な解析的・計算的性質を持ち、\emph{エントロピー関数}に基づいている。

\begin{defn}[エントロピー関数]
\label{def_entropy}
関数$\phi : \RR \to \RR \cup \{\infty\}$がエントロピー関数であるとは、下半連続かつ凸で、$\dom \phi\subset [0,\infty[$であり、以下の実行可能性条件$\dom \phi \; \cap\;  ]0, \infty[\; \neq \emptyset$を満たすことをいう。$\phi$の無限遠での増大の速度は以下で記述される。%its recession (or horizon) function
\eq{
\phi'_\infty = \lim_{x\rightarrow +\infty} \varphi(x)/x \in \RR \cup \{\infty\} \, .
}
\end{defn}

$\phi'_\infty = \infty$のとき、$\phi$は任意の線形関数より速く増大し、$\phi$は\emph{超線形}であるという。任意のエントロピー関数$\phi$は、以下のように$\phi$-ダイバージェンス（Cisz\'arダイバージェンス~\citep{ciszar1967information,ali1966general}または$f$-ダイバージェンスとも呼ばれる）を誘導する。



\begin{defn}[$\phi$-ダイバージェンス]
\label{def_divergence}
$\phi$をエントロピー関数とする。
$\al,\be \in \Mm(\X)$に対し、$\frac{\d \al}{\d \be} \be + \al^{\perp}$を$\be$に関する$\al$のLebesgue分解\footnote{Lebesgueの分解定理は、$\be$が与えられたとき、$\al$は二つの測度$\al^s + \al^\perp$の和として一意に分解され、$\al^s$は$\be$に関して絶対連続であり、$\al^\perp$と$\be$は互いに特異であることを主張する。}とする。ダイバージェンス$\Divergm_\phi$は以下で定義される。
\eql{\label{eq-phi-div}
	\Divergm_\phi (\al|\be) \eqdef \int_\X \phi\left(\frac{\d \al}{\d \be} \right) \d \be
+ \phi'_\infty \al^{\perp}(\X)
}
ただし$\al,\be$が非負の場合であり、そうでなければ$\infty$とする。
\end{defn}%

~\eqref{eq-phi-div}における追加項$\phi'_\infty \al^{\perp}(\X)$は、$\phi$が無限遠で線形増大を持つ場合でも（例えば、TVノルムを定義する絶対値~\eqref{eq-tv-entropy}の場合のように）、$\Divergm_\phi$が（測度の弱位相に関する）連続汎関数を定義することを保証するために重要である。$\phi$が超線形増大を持つ場合、例えば通常のエントロピー~\eqref{eq-shannon-entropy}では、$\phi'_\infty=+\infty$となるため、$\al$が$\be$に関して密度を持たなければ$\Divergm_\phi (\al|\be) = +\infty$となる。

離散的な設定では、
\eql{\label{eq-div-disc-meas}
	\al=\sum_i \a_i \de_{x_i}
	\qandq \be=\sum_i \b_i \de_{x_i}
}
が$n$個の点$(x_i)_{i=1}^n \subset \X$の同一集合上に台を持つと仮定すると、~\eqref{eq-phi-div}は$\simplex_n$上のダイバージェンスを定義する。
\eql{\label{eq-discr-diverg}
	\DivergmD_\phi(\a|\b) = \sum_{i \in \Supp(\b)} \phi\pa{ \frac{\a_i}{\b_i} } \b_i + \phi'_\infty \sum_{i \notin \Supp(\b)} \a_i,
}
ここで$\Supp(\b) \eqdef \enscond{i \in \range{n}}{ b_i \neq 0 }$である。

以下の命題の証明は~\cite[Thm 2.7]{LieroMielkeSavareLong}に見ることができる。

\begin{proposition}
$\phi$がエントロピー関数であるならば、$\Divergm_\phi$は$(\al,\be)$について同次1次、凸、かつ弱*下半連続である。
\end{proposition}


\begin{figure}[h!]
\centering
\includegraphics[width=.6\linewidth]{divergences/divergences}
\caption{\label{fig-divergences}
エントロピー汎関数の例。
}
\end{figure}


\begin{rem}[双対表現]
	$\phi$-ダイバージェンスは、$\phi$のLegendre変換
	\eq{
		\phi^*(s) \eqdef \usup{t \in \RR} st - \phi(t)
	}
	（~\eqref{eq-legendre}も参照）を用いて以下のように表現できる。
	\eq{
		\Divergm_\phi (\al|\be) = \usup{f: \X \rightarrow \RR} \int_\X f(x) \d\al(x) - \int_\X \phi^*(f(x)) \d\be(x);
	}
	詳細は~\citet{LieroMielkeSavareLong}を参照されたい。
\end{rem}


ここで、この枠組みのいくつかの代表的な例を概観する。図~\ref{fig-divergences}は対応するエントロピー汎関数を示し、図~\ref{fig-divergences-relations}はそれらの関係を概観する。





\begin{figure}[h!]
\centering
\includegraphics[width=.6\linewidth]{divergences/divergences-comparison}
\caption{\label{fig-divergences-relations}
ダイバージェンス間の関係図（~\citet{gibbs2002choosing}に着想を得たもの）。
$\X$を基底距離$d$を持つ距離空間とし、$d_{\max} = \sup_{(x,x')} d(x,x')$を$\X$の直径とする。$\X$が離散的な場合、$d_{\min} \eqdef \min_{x \neq x'} d(x,x')$とする。\todoK{Tripple check these bounds since they might use different definition/normalization. }\todoK{Change the notation for the Hellinger distance one the symbol is decided $\Hellinger$. }
}
\end{figure}


\begin{example}[Kullback--Leiblerダイバージェンス]
\label{ex_KLdiv}
Kullback--Leiblerダイバージェンス$\KL \eqdef \Divergm_{\phi_{\KL}}$は、相対エントロピーとも呼ばれ、既に~\eqref{eq-defn-rel-entropy}と~\eqref{eq-kl-defn}で導入された。これはShannon--Boltzmannエントロピー関数$\phi_{\KL}$に対応するダイバージェンスであり、以下で与えられる。
\eql{\label{eq-shannon-entropy}
	\phi_{\KL}(s)= \begin{cases}
		s\log(s)-s+1 & \textnormal{for } s>0 , \\
		1 & \textnormal{for } s=0 , \\
		+\infty & \textnormal{otherwise.}
		\end{cases}
}
\end{example}

\begin{remark}[Bregmanダイバージェンス]\label{rem-bregman}
離散KLダイバージェンス$\KLD \eqdef \DivergmD_{\phi_{\KL}}$は、$\phi$-ダイバージェンスでありかつBregmanダイバージェンスでもあるという一意的な性質を持つ。
%
$\RR^n$上の離散ベクトルに対して、滑らかな狭義凸関数$\psi : \RR^n \rightarrow \RR$に対応するBregmanダイバージェンス~\citep{bregman1967relaxation}は以下で定義される。
\eql{\label{eq-bregman-discr}
	\VectMode{B}_\psi(\a|\b) \eqdef \psi(\a) - \psi(\b) - \dotp{\nabla \psi(\b)}{\a-\b},
}
ここで$\dotp{\cdot}{\cdot}$は$\RR^n$上の標準内積である。$\VectMode{B}_\psi(\a|\b)$は$\a$の凸関数であり、$\psi$の線形関数であることに注意されたい。
%
$\phi$-ダイバージェンスと同様に、Bregmanダイバージェンスは$\VectMode{B}_\psi(\a|\b) \geq 0$を満たし、$\VectMode{B}_\psi(\a|\b)=0$であることと$\a=\b$であることは同値である。
%
KLダイバージェンスは、~\eqref{eq-discr-entropy}で定義された負のエントロピー$\psi = -\HD$に対するBregmanダイバージェンスであり、すなわち$\KLD = \VectMode{B}_{-\HD}$である。
%
Bregmanダイバージェンスは局所的にはEuclid距離の二乗であり、
\eq{
	\VectMode{B}_\psi(\a + \varepsilon|\a + \eta) = \dotp{ \partial^2 \psi(\a)(\varepsilon-\eta) }{\varepsilon-\eta}
	+ o(\norm{\varepsilon-\eta}^2)
}
が成り立ち、分離点の集合$\enscond{\a}{\VectMode{B}_\psi(\a|\b) = \VectMode{B}_\psi(\a|\b')}$は$\b$と$\b'$の間の超平面である。
%
これらの性質により、Bregmanダイバージェンスは一次最適化法においてEuclid距離の代替として適している。最も有名な例はミラー勾配降下法~\citep{beck2003mirror}であり、これは~\eqref{eq-explicit-metricflow}の形式の陽的降下ステップである。
%
Bregmanダイバージェンスは凸最適化においても重要であり、例えばSinkhorn反復の導出や有限次元での収束の研究に用いることができる。注釈~\ref{rem-iterative-projection}を参照されたい。
%%%% Should be re-written in infinite dimension
%While KL is both a $\phi$ and a Bregman divergence, these are radically different concepts. In particular, to generalize formula~\eqref{eq-bregman-discr} to nondiscrete measures on some space $\Xx$, one needs to fix a reference measure $\xi$ and consider only measures $\al,\be$ with square integrable densities $\frac{\d\al}{\d\xi},\frac{\d\be}{\d\xi} \in L^2(\d\xi)$. For some strictly convex $\psi : L^2(\d\xi) \rightarrow \RR$, one then defines the Bregman divergence as
%\eq{
%	\Bb_{\psi,\xi}(\al|\be) \eqdef \psi\pa{\frac{\d\al}{\d\xi}} -
%		\psi\pa{\frac{\d\be}{\d\xi}} -
%		\int_{\X} \eta(x) \pa{
%			\frac{\d\al}{\d\xi}(x)-\frac{\d\be}{\d\xi}(x)
%		} \d \xi(x),
%}
%where $\eta = \nabla \psi\pa{\frac{\d\be}{\d\xi}}  \in L^2(\d\xi)$ is the gradient~\citep{rao1985cross,jones1990general}.
%%%
%The simplest case is if $\psi(\rho) = \int_\X \tilde\psi(\rho(x)) \d\xi(x)$ in which case
%\eq{
%	\Bb_{\psi,\xi}(\al|\be) \eqdef
%		\int_\X \pa{
%		\tilde\psi\pa{\frac{\d\al}{\d\xi}} -
%		\tilde\psi\pa{\frac{\d\be}{\d\xi}} -
%		\tilde\psi'\pa{\frac{\d\be}{\d\xi}}
%		\pa{
%			\frac{\d\al}{\d\xi}-\frac{\d\be}{\d\xi}
%		}
%		} \d \xi.
%}
\end{remark}

\begin{remark}[KLの双曲幾何]
	Kullback--Leiblerダイバージェンスの幾何学を、Gauss分布の比較における二次最適輸送が定義する幾何学と対比することは興味深い。例えば~\citet{costa2015fisher}が詳述しているように、Kullback--LeiblerダイバージェンスはGauss密度に対して閉じた形を持つ。一変量の場合$\dim=1$において、$\al = \Nn(m_\al,\si_\al^2)$かつ$\be = \Nn(m_\be,\si_\be^2)$とすると、
	\eql{\label{eq-kl-gaussian}
		\KL(\al|\be) = \frac{1}{2}\pa{
			\frac{\si_\al^2}{\si_\be^2}  +
			\log\pa{\frac{\si_\be^2}{\si_\al^2}}
			+ \frac{|m_\al-m_\be|}{\si_\be^2}
			 - 1
		}.
	}
	この表現は、$\si_\be$が$0$に減少し$\be$がDirac質量になるにつれて$\al$と$\be$の間のダイバージェンスが無限大に発散することを示している。この意味で、特異Gauss分布はKL幾何において他のすべてのGauss分布から無限に遠いと言える。したがって、この幾何学は特異共分散を扱うことを避けたい場合に有用である。分析を簡略化するために、KLの無限小幾何学を見ることができ、これは2次のテイラー展開によって得られる。
	\eq{
		\KL( \Nn(m+\de_m,(\si+\de_\si)^2) | \Nn(m,\si^2) ) = \frac{1}{\si^2}\pa{
			\frac{1}{2}\de_m^2 + \de_\si^2
		} +
		o( \de_m^2,\de_\si^2 ).
	}
	この局所Riemann計量、いわゆるFisher計量は、$(m/\sqrt{2},\si) \in \RR \times \RR_{+,*}$上で表現すると、双曲Poincaré半平面の計量と正確に一致する。
	%
	この空間上の測地線は$\si=0$の直線を中心とする半円であり、指数的な速度を持つ。すなわち、極限$\si=0$に到達するのは無限時間後である。
	%
	特に、$\si_\al=\si_\be$であるが$m_\al \neq m_\al$の場合、この双曲半平面上の$(\al,\be)$間の測地線は一定の標準偏差を持たないことに注意されたい。

	Gaussパラメータ$(m,\si)$の空間上のKL双曲幾何学は、注釈~\ref{rem-dist-gaussians}で記述されたOTに付随するEuclid幾何学と対比されるべきである。一変量の場合、
	\eql{\label{eq-univariate-gauss}
		\Wass_2^2( \al,\be ) = |m_\al - m_\be |^2 + |\si_\al-\si_\be|^2.
	}
	図~\ref{fig-kl-gaussians}は、これら二つの幾何学とそれぞれの測地線の視覚的比較を示している。
	%
	この興味深い比較はJean Feydyによって示唆されたものである。
\end{remark}



\begin{figure}[h!]
\centering
\begin{tabular}{@{}c@{\hspace{3mm}}c@{}}
\includegraphics[width=.45\linewidth]{kl-gaussian/interp-kl}&
\includegraphics[width=.46\linewidth]{kl-gaussian/interp-ot}\\
KL & OT
\end{tabular}
\caption{\label{fig-kl-gaussians}
KL（双曲）幾何とOT（Euclid）幾何を用いたGauss分布間の補間の比較。
}
\end{figure}


\begin{example}[全変動]\label{exmp-tv}
全変動距離$\TV \eqdef \Divergm_{\phi_{\TV}}$は、以下に対応するダイバージェンスである。
\eql{\label{eq-tv-entropy}
	\phi_{\TV}(s)= \begin{cases}
		|s-1| & \textnormal{for } s\geq0 , \\
		+\infty & \textnormal{otherwise.}
		\end{cases}
}
これは実際には測度の全空間$\Mm(\X)$上のノルムを定義し、
\eql{\label{eq-defn-tv}
	\TV(\al|\be) = \norm{\al-\be}_{\TV},
	\qwhereq
	\norm{\al}_{\TV} = |\al|(\X) = \int_\X \d|\al|(x).
}
となる。$\al$が$\X=\RR^\dim$上で密度$\density{\al}$を持つ場合、TVノルムは関数の$L^1$ノルムであり、$\norm{\al}_{\TV} = \int_\X |\density{\al}(x)| \d x = \norm{\density{\al}}_{L^1}$となる。
%
$\al$が~\eqref{eq-div-disc-meas}のように離散的である場合、TVノルムは$\RR^n$のベクトルの$\ell^1$ノルムであり、$\norm{\al}_{\TV}=\sum_i |\a_i| = \norm{\a}_{\ell^1}$となる。
\end{example}

\begin{rem}[強位相と弱位相]
	全変動ノルム~\eqref{eq-defn-tv}は測度の空間上のいわゆる「強」位相を定義する。
	%
	半径$R$のコンパクト領域$\X$上では、
	\eq{
		\Wass_1(\al,\be) \leq R \norm{\al-\be}_{\TV}
	}
	が成り立つので、この強い収束の概念はWasserstein距離によって距離化される弱収束を含意する。
	%
	しかし逆は成り立たない。$x \rightarrow y$のとき$\de_x$は$\de_y$に強収束しないからである（$x \neq y$のとき$\norm{\de_x-\de_y}_{\TV}=2$であることに注意）。
	%
	主要な利点は、$\Mm_+^1(\Xx)$が（再びコンパクトな台空間$\X$上で）弱位相に関してコンパクトであることであり、したがって確率測度の任意の列$(\al_k)_k$から常に収束する部分列を取り出すことができる。これにより、第~\ref{c-variational}章で考察するような最適化問題に適した空間となる。
\end{rem}


\begin{example}[Hellinger]\label{exmp-hellinger}
	Hellinger距離$\Hellinger \eqdef \Divergm_{\phi_{H}}^{1/2}$は、以下に対応するダイバージェンスの平方根である。
	\eq{
		\phi_{H}(s)= \begin{cases}
			|\sqrt{s}-1|^2 & \textnormal{for } s\geq0 , \\
			+\infty & \textnormal{otherwise.}
			\end{cases}
	}
	名前が示す通り、$\Hellinger$は$\Mm_+(\X)$上の距離であり、$\norm{\cdot}_{\TV}$と同様に強位相を距離化する。
	%
	$(\al,\be)$が$\X=\RR^\dim$上で密度$(\density{\al},\density{\be})$を持つ場合、$\Hellinger(\al,\be) = \norms{\sqrt{\density{\al}}-\sqrt{\density{\be}}}_{L^2}$となる。
	%
	$(\al,\be)$が~\eqref{eq-div-disc-meas}のように離散的である場合、$\Hellinger(\al,\be) = \norms{\sqrt{\a}-\sqrt{\b}}$となる。
	%
	$\phi_{L^p}(s)=|s^{1/p}-1|^p$を考えるとHellinger（$p=2$）および全変動（$p=1$）距離が一般化され、$\Divergm_{\phi_{L^p}}^{1/p}$は$0<p<+\infty$に対して強収束を距離化する距離となる。
\end{example}


\begin{example}[Jensen--Shannon距離]
	KLダイバージェンスは対称ではなく、Bregmanダイバージェンス（局所的に二次ノルム）ではあるが、距離の二乗ではない。一方、Jensen--Shannon距離$\text{JS}(\al,\be)$は以下で定義され、
	\eq{
		\text{JS}(\al,\be)^2 \eqdef \frac{1}{2}\pa{
			\KL(\al|\xi) + \KL(\be|\xi)
		}
		\qwhereq
		\xi = \frac{\al+\be}{2},
	}
	距離である~\citep{endres2003new,osterreicher2003new}。
	%
	$\text{JS}^2$は$\phi(s) = t\log(t)-(t+1)\log(t+1)$に対する$\phi$-ダイバージェンスであることが示せる。
	%
	$\KL$とは鋭く対照的に、$\text{JS}(\al,\be)$は常に有界である。より正確には、$0 \leq \text{JS}(\al,\be)^2 \leq \ln(2)$を満たす。
	%
	TVノルムやHellinger距離と同様に、強収束を距離化する。
\end{example}



\begin{example}[$\chi^2$]\label{exmp-chisquare}
	$\chi^2$-ダイバージェンス$\chi^2 \eqdef \Divergm_{\phi_{\chi^2}}$は、以下に対応するダイバージェンスである。
	\eq{
		\phi_{\chi^2}(s)= \begin{cases}
			|s-1|^2 & \textnormal{for } s\geq0 , \\
			+\infty & \textnormal{otherwise.}
			\end{cases}
	}
	$(\al,\be)$が~\eqref{eq-div-disc-meas}のように離散的で同一の台を持つ場合、
	\eq{
		\chi^2(\al|\be) = \sum_i \frac{(\a_i-\b_i)^2}{\b_i}.
	}
\end{example}


%%%%%%%%%%%%%%%%%%%%%%%%%%%%%%%%%%%%%%%%%%%%%%%%%%%%%%%%%%%%%%%%%%%%%%%%%%%%%%%%%%%%%%%%%%%%%%%%%%%%
%%%%%%%%%%%%%%%%%%%%%%%%%%%%%%%%%%%%%%%%%%%%%%%%%%%%%%%%%%%%%%%%%%%%%%%%%%%%%%%%%%%%%%%%%%%%%%%%%%%%
%%%%%%%%%%%%%%%%%%%%%%%%%%%%%%%%%%%%%%%%%%%%%%%%%%%%%%%%%%%%%%%%%%%%%%%%%%%%%%%%%%%%%%%%%%%%%%%%%%%%
\section{積分確率計量}
\label{sec-dual-norms}


定式化~\eqref{eq-w1-cont}は双対ノルムの特殊な場合である。双対ノルムは、任意の測度を扱うことのできる「弱い」ノルムを設計するための便利な方法である。可測関数の対称凸集合$B$に対して、以下を定義する。
\eql{\label{eq-dual-norm-cont}
	\norm{\al}_B \eqdef
	\umax{\f} \enscond{ \int_{\X} \f(x) \d\al(x) }{ \f \in B}.
}
これらの双対ノルムは「積分確率計量」と呼ばれることが多い。~\citep{sriperumbudur2012empirical}を参照されたい。


% Cf. ``Integral probability metrics and their generating classes of functions".

\begin{example}[全変動]
全変動ノルム（例~\ref{exmp-tv}）は、連続関数の全空間に対応する双対ノルムである。
\eq{
	B = \enscond{f \in \Cc(\X)}{\norm{f}_\infty \leq 1}.
}
全変動距離は、双対ノルムでもある唯一の非自明なダイバージェンスである。~\citep{sriperumbudur2009integral}を参照されたい。
\end{example}

\begin{rem}[弱収束の距離化]\label{rem-ipm-weak}
より小さな「球」$B$を用いることで、典型的には連続（ときに正則な）関数のみを含む、より弱い双対ノルムが定義される。
%
$\norm{\cdot}_B$が弱収束（定義~\ref{dfn-weak-conv}を参照）を距離化するためには、$B$で張られる空間がsupノルム$\norm{\cdot}_\infty$（すなわち一様収束の位相）に関して連続関数の集合の中で稠密であれば十分である。~\citep[para. 5.1]{ambrosio2006gradient}を参照されたい。
\end{rem}

図~\ref{fig-dual-norms}はいくつかのこのような双対ノルムの比較を示しており、以下で詳述する。


\begin{figure}[h!]
\centering
\begin{tabular}{@{}c@{\hspace{3mm}}c@{}}
\includegraphics[width=.45\linewidth]{dual-norms/1dirac-1dirac} &
\includegraphics[width=.45\linewidth]{dual-norms/1dirac-2dirac} \\
$(\al,\be)=(\de_0,\de_t)$ &
$(\al,\be)=(\de_0,\frac{1}{2}(\de_{-t/2}+\de_{t/2}))$
\end{tabular}
\caption{\label{fig-dual-norms}
双対ノルムの比較。
}
\end{figure}



%%%%
\subsection{$\Wass_1$とフラットノルム}

集合$B$が有界であれば、$\norm{\cdot}_B$は測度の全空間$\Mm(\Xx)$上のノルムとなる。
%
$\Wass_1$の場合はそうではなく、$\int_\X \d\al=0$を満たす$\al$に対してのみ定義される（そうでなければ$\norm{\al}_B=+\infty$となる）。
%
これはポテンシャル$\f$の値に上界を課すことで緩和でき、例えばフラットノルムを定義する。

\begin{example}[$\Wass_1$ノルム]
~\eqref{eq-w1-cont}で定義される$\Wass_1$は、以下を用いる双対ノルム~\eqref{eq-dual-norm-cont}の特殊な場合である。
\eq{
	B = \enscond{f}{\Lip(f) \leq 1}
}
すなわち1-Lipschitz関数の集合である。
\end{example}

\begin{example}[フラットノルムとDudley計量]
フラットノルムは以下を用いて定義される。
\eql{\label{eq-set-flatnorm}
	B=\enscond{f}{\norm{\nabla \f}_\infty \leq 1 \qandq \norm{\f}_\infty \leq 1}.
}
これは測度の全空間$\Mm(\X)$上で弱収束を距離化する。
%
式~\eqref{eq-w1-discr}は、制約$\abs{\fD_k} \leq 1$を追加することでフラットノルムの計算に拡張される。
%
フラットノルムは「Kantorovich--Rubinstein」ノルム~\citep{hanin1992kantorovich}と呼ばれることもあり、画像処理における逆問題の忠実度項として用いられてきた~\citep{lellmann2014imaging}。
%
フラットノルムはDudley計量と類似しており、こちらは以下を用いる。
\eq{\label{eq-set-dudley}
	B=\enscond{f}{\norm{\nabla \f}_\infty + \norm{\f}_\infty \leq 1}.
}
\end{example}


%%%%
\subsection{双対RKHSノルムと最大平均差異}
\label{sec-mmd}


カーネル法、より具体的には再生核Hilbert空間（RKHS; データサイエンスにおける応用の概説は~\citep{scholkopf2002learning}を参照）の機構を用いて、測度上の「Euclid的」ノルム（二次汎関数を用いて構成される）を定義することも可能である。まず基本的な定義を想起する。

\begin{defn}\label{def-negativedefinitekernel}
集合$\X \times \X$上で定義された対称関数$k$（あるいは$\varphi$）が正定値（あるいは負定値）であるとは、任意の$n\geq0$、任意の族$x_1,\dots,x_n\in\Z$、および任意のベクトル$r\in\RR^n$に対して以下の不等式が成り立つことをいう。
\eql{\label{eq-dual-kern}
	\sum_{i,j=1}^n r_i r_j k(x_i,x_j)\geq 0,
	\quad
	\left(\text{resp.}\quad\sum_{i,j=1}^n r_i r_j \varphi(x_i,x_j)\leq 0\right).
}
カーネルが条件付き正定値であるとは、~\eqref{eq-dual-kern}において正定値性が零平均ベクトル$r$（すなわち$\dotp{r}{\ones_n}=0$を満たすもの）に対してのみ成り立つ場合をいう。
\end{defn}


\if 0 %%%% REMOVED
Using a positive definite kernel $k$, one can now define a Hilbertian norm $\norm{\cdot}_\Krkhs$ on functions
\eq{
	\norm{f}_\Krkhs^2 \eqdef \int_{\X \times \X} \Krkhs(x,y) f(x) f(y) \d \rho(x) \d \rho(y)
}
where $\rho$ is some reference measure, typically Lebesgue's measure on $\X=\RR^\dim$. This norm can then be used to define a subset $B$ of functions with bounded RKHS norm,
\eql{\label{eq-dual-rkhs}
	B = \enscond{\f}{ \norm{\f}_\Krkhs \leq 1}
}
%
that is plugged into notation~\eqref{eq-dual-norm-cont}. The dual RKHS norm for measures that results from this construction is often referred to as the ``maximum mean discrepancy'' (MMD) (see~\citet{gretton2007kernel}), and can be revisited through the prism of Kernel Mean Embeddings (see~\citet{muandet2017kernel} for a review).
\fi


$k$が条件付き正定値のとき、以下のノルムが定義される。
\eql{\label{eq-kernel-dual}
	\norm{\al}^2_{\Krkhs} \eqdef \int_{\X \times \X} \Krkhs(x,y) \d \al(x) \d \al(y).
}
これらのノルムは「最大平均差異」（MMD）と呼ばれることが多く（~\citep{gretton2007kernel}を参照）、形状解析では「カーネルノルム」とも呼ばれてきた~\citep{glaunes2004diffeomorphic}。
%
この表現~\eqref{eq-kernel-dual}は、$\X$上の分布$\al$に従う二つの独立な確率ベクトル$(X,X')$を導入して以下のように言い換えることができる。
\eq{
	\norm{\al}^2_{\Krkhs} = \EE_{X,X'}( \Krkhs(X,X') ).
}
$\norm{\cdot}^2_{\Krkhs}$は、$k$に対応するRKHSの単位球$B$に対応する~\eqref{eq-dual-norm-cont}の意味での双対ノルムであることが示される。RKHS関数空間の詳細は~\citep{berlinet03reproducing,Hofmann2008,scholkopf2002learning}を参照されたい。

\begin{rem}[普遍カーネル]
注釈~\ref{rem-ipm-weak}によれば、MMDノルム$\norm{\cdot}_k$は、双対球$B$の張る空間が連続関数の空間$\Cc(\Xx)$の中で稠密であれば弱収束を距離化する。これは、有限和$\sum_{i=1}^n a_i k(x_i,\cdot)$（$n$と点$(x_i)_i$の任意の選択に対して）が一様ノルム$\norm{\cdot}_\infty$に関して$\Cc(\Xx)$の中で稠密であることを意味する。$\X=\RR^d$上の平行移動不変カーネル$\Krkhs(x,y)=\Krkhs_0(x-y)$の場合、これはFourier変換が消えないこと$\hat \Krkhs_0(\om) >0$と同値である。
\end{rem}

$\al$が~\eqref{eq-pair-discr}の形の離散測度である特殊な場合には、以下の簡単な表現が得られる。
\eq{
	\norm{\al}_{\Krkhs}^2 = \sum_{i=1}^n \sum_{i'=1}^n \a_i\a_{i'} \KrkhsD_{i,i'}  = \dotp{\KrkhsD\a}{\a}
	\qwhereq
	\KrkhsD_{i,i'} \eqdef \Krkhs(x_i,x_{i'}).
}
特に、$\al=\sum_{i=1}^n \a_i \de_{x_i}$と$\be=\sum_{i=1}^n \b_i \de_{x_i}$が同一の点集合上に台を持つとき、$\norm{\al-\be}_{\Krkhs}^2 = \dotp{\KrkhsD(\a-\b)}{\a-\b}$となるので、$\norm{\cdot}_{\Krkhs}$は単体$\simplex_n$上のEuclidノルム（$\KrkhsD$が正定値なら真のノルム、$\KrkhsD$が半正定値なら退化したノルム）である。
%
~\eqref{eq-pair-discr}の形の二つの離散測度間の差異を計算するには、以下を用いることができる。
\eql{\label{eq-mmd-discr}
	\norm{\al-\be}_{\Krkhs}^2 =
		\sum_{i,i'} \a_i \a_{i'} \Krkhs(x_i,x_{i'})	+
		\sum_{j,j'} \b_j \b_{j'} \Krkhs(y_j,y_{j'}) - 2
		\sum_{i,j} \a_i \b_j \Krkhs(x_i,y_j).
}

\begin{example}[GaussRKHS]\label{exmp-gaussian-kernel}
最もよく用いられるカーネルの一つはGaussカーネル$\Krkhs(x,y) = e^{-\frac{\norm{x-y}^2}{2 \si^2}}$であり、これは$\Xx=\RR^d$上の正定値普遍カーネルである。
%
Gaussカーネルの魅力的な特徴は、1次元カーネルの積として分離可能であることであり、これにより正則格子上での計算が容易になる（注釈~\ref{rem-separable}も参照）。
しかし、Gaussカーネルを用いる際に生じる重要な問題は、帯域幅パラメータ$\si$を選択する必要があることである。この帯域幅は、比較する測度の観測値間の「典型的なスケール」と一致する必要がある。測度がマルチスケールの特徴を持つ場合（ある領域は非常に密で、他の領域は非常に疎である場合）、Gaussカーネルは適切ではなく、次に詳述する「スケールフリー」カーネルを検討すべきである。このようなスケールフリーカーネルの問題は、大域的（遅い多項式減衰を持つ）であることであり、コンパクト台近似が不可能なため、計算コストが典型的に高くなる。
%
図~\ref{fig-rkhs}はいくつかのカーネルの比較を示している。
\end{example}

\begin{figure}[h!]
\centering
\begin{tabular}{@{}c@{\hspace{0mm}}c@{\hspace{0mm}}c@{\hspace{0mm}}c@{\hspace{0mm}}c@{}}
 &
\myFigRKHS{energy-dist-kernel}&
\myFigRKHS{gaussian-small-kernel}&
\myFigRKHS{gaussian-medium-kernel}&
\myFigRKHS{gaussian-large-kernel}\\
\myFigRKHS{input} &
\myFigRKHS{energy-dist} &
\myFigRKHS{gaussian-small}&
\myFigRKHS{gaussian-medium}&
\myFigRKHS{gaussian-large}\\
$({\color{red} \al},{\color{blue} \be})$ & $\text{ED}(\RR^2,\norm{\cdot})$ & $(G,.005)$ &  $(G,.02)$ & $(G,.05)$
\end{tabular}
\caption{\label{fig-rkhs}
上段: $\norm{\al-\be}_k = \norm{\psi \star (\al-\be)}_{L^2(\RR^2)}$なる$\psi$の表示。Fourier上では$\hat\psi(\om) = \sqrt{ \hat k_0(\om) }$として形式的に定義される。ここで$\Krkhs(x,x')=k_0(x-x')$である。
%
下段: $\psi \star (\al-\be)$の表示。
%
(G,$\si$)は分散$\si^2$のGaussカーネルを表す。$\text{ED}(\RR^2,\norm{\cdot})$のカーネルは$\psi(x)=1/\sqrt{\norm{x}}$である。
%
\todoK{check that in 2D the ED sqrt kernel is $1/\sqrt{|x|}$, I think this is the case.} \todoK{Explain that ED and large sigma are too global to match the feature size, medium sigma only capture one scale of the details ; small sigma too small}
}
\end{figure}


\begin{example}[$H^{-1}(\RR^\dim)$]\label{exp-sobolev-neg1}
もう一つの重要な双対ノルムは$H^{-1}(\RR^\dim)$であり、これは$L^2(\RR^\dim)$に導関数を持つ関数のSobolev空間$H^1(\RR^\dim)$の（超関数上の）双対である。
%
これは主RKHSノルム$\norm{\nabla \f}_{L^2(\RR^\dim)}^2$を用いて定義される。$\dim=1$でない限り、$H^1(\RR^\dim)$のSobolev空間の関数は一般に連続でないため、特異測度（例えばDirac）に対しては定義されない。
%
この$H^{-1}$ノルム（零平均で密度を持つ測度の空間上で定義される）はダイバージェンス形式でも定式化できる。
\eql{\label{eq-dual-sobolev-div}
	\norm{\al-\be}_{H^{-1}(\RR^\dim)}^2 =
	\umin{\flow} \enscond{ \int_{\RR^\dim} \norm{\flow(x)}_2^2 \d x }{  \diverg(\flow)=\al-\be },
}
これは~\eqref{eq-w1-cont-div}と対比されるべきであり、そちらではここで用いた$L^2$ノルムの代わりにベクトル場$\flow$の$L^1$ノルムが用いられていた。
%
このSobolev双対ノルムの「重み付き」版、
\eq{
	\norm{\rho}_{H^{-1}(\al)}^2 =
	\umin{\diverg(\flow)=\rho} \int_{\RR^\dim} \norm{\flow(x)}_2^2 \d \al(x),
}
はWasserstein $\Wass_2$ノルムの自然な「線形化」として解釈できる。すなわち、Benamou--Brenier動的定式化は無限小的に以下のように解釈できる。
\eql{\label{eq-wass-asymp-sob}
	\Wass_2(\al,\al+\varepsilon \rho) = \varepsilon \norm{\rho}_{H^{-1}(\al)} + o(\varepsilon).
}
汎関数$\Wass_2(\al,\be)$と$\norm{\al-\be}_{H^{-1}(\al)}$は同値であることが示される~\citep{peyre2011comparison}。
%
問題は、$\norm{\al-\be}_{H^{-1}(\al)}$がノルムではないこと（$\al$による重み付けのため）であり、$(\al,\be)$が密度を持たない限り、一般にはこれを$\norm{\al-\be}_{H^{-1}(\RR^\dim)}$で置き換えることはできない。
%
この場合、$\al$と$\be$が同一の台上で下から$a>0$、上から$b<+\infty$で有界な密度を持つならば、
\eql{\label{eq-compar-w2-sob}
	b^{-1/2} \norm{\al-\be}_{H^{-1}(\RR^\dim)} \leq W_2(\al,\be) \leq a^{-1/2} \norm{\al-\be}_{H^{-1}(\RR^\dim)};
}
が成り立つ。~\citep[Theo. 5.34]{SantambrogioBook}を参照し、最適定数については~\citep{peyre2011comparison}を参照されたい。
%
% Note however $H^{-1}(\al)$ and $\Wass_2(\al,\be)$ are related when comparing close enough distribution, they have very different behaviors.
\end{example}

\begin{example}[負Sobolev空間]\label{exp-sobolev-neg}
	この構成を一般化するには、任意の負の指数のSobolev空間$H^{-r}(\RR^\dim)$を考えればよい。これは$r$階導関数（超関数の意味で）が$L^2(\RR^\dim)$に属する関数のSobolev空間$H^{r}(\RR^\dim)$の双対である。
	%
	弱収束を距離化するためには、$H^{r}(\RR^\dim)$の関数が連続である必要があり、これは$r>\dim/2$のとき成り立つ。次元$\dim$が増加すると、より高い正則性を考える必要がある。
	%
	任意の$\al$（必ずしも整数でなくてよい）に対して、これらの空間はFourier変換を用いて定義され、Fourier変換$\hat \al(\om)$（ここではLebesgue測度$\d\om$に関する密度として記述）を持つ測度$\al$に対して、
	\eq{
		\norm{\al}_{H^{-r}(\RR^\dim)}^2 \eqdef \int_{\RR^\dim} \norm{\om}^{-2r} |\hat \al(\om)|^{2}\d\om.
	}
	これは$\hat k_0(\om) = \pm \norm{\om}^{-2r}$を持つ畳み込みカーネル$k(x,y)=k_0(x-y)$を用いた双対RKHSノルムに対応する。逆Fourier変換を取ると、（定数を除いて）以下が得られることがわかる。
	 \eql{\label{eq-sob-kern}
	 	\foralls x \in \RR^\dim, \quad
		k_0(x) =
		\choice{
			\frac{1}{\norm{x}^{d-2r}} \qifq r < d/2,  \\
			-\norm{x}^{2r-d} \qifq r>d/2.
		}
	}
\end{example}




\begin{example}[エネルギー距離]\label{exp-energy-dist}
エネルギー距離（$\dim=1$のときCramer距離）~\citep{szekely2004testing}は、距離$\dist$に対して以下で定義される。
\eql{\label{eq-defn-ed}
	\norm{\al-\be}_{\text{ED}(\X,\dist^p)} \eqdef \norm{\al-\be}_{\Krkhs_{\text{ED}}}
	\qwhereq
	\Krkhs_{\text{ED}}(x,y) = -\dist(x,y)^p
}
ただし$0 < p < 2$とする。\todoK{is it correct?}
%
$\dist$が負定値（定義~\ref{def-negativedefinitekernel}を参照）であれば、これは測度上の有効なMMDノルムであり、典型的な例はEuclid距離$\dist(x,y)=\norm{x-y}$である。
%
$\X=\RR^\dim$、$\dist(x,y)=\norm{\cdot}$の場合、~\eqref{eq-sob-kern}を用いると、エネルギー距離はSobolevノルムであることがわかる。
\eq{
	\norm{\cdot}_{\text{ED}(\RR^\dim,\norm{\cdot}^p)} = \norm{\cdot}_{H^{-\frac{d+p}{2}}(\RR^\dim)}.
}
エネルギー距離がGaussカーネル（例~\ref{exmp-gaussian-kernel}）のようなより一般的なカーネルに対して持つ主要な利点は、スケールフリーであり帯域幅パラメータ$\si$に依存しないことである。より正確には、$\X=\RR^\dim$上で因子$s>0$による膨張を$f_s(x)=sx$と記すと、以下のスケーリング挙動を持つ。
\eq{
	\norm{f_{s\sharp}(\al-\be)}_{\text{ED}(\RR^\dim,\norm{\cdot}^p)} = s^{\frac{p}{2}} \norm{\al-\be}_{\text{ED}(\RR^\dim,\norm{\cdot}^p)},
}
一方、Wasserstein距離は完全な線形スケーリングを示す。
\eq{
	\Wass_p( f_{s\sharp}\al,f_{s\sharp}\be) ) = s \Wass_p(\al,\be) ).
}
ただし、エネルギー距離ではパラメータ$p$は$0 < p < 2$を満たす必要があり、$p=2$のときは平均間の距離に退化する。
\eq{
	\norm{\al-\be}_{\text{ED}(\RR^\dim,\norm{\cdot}^2)}
	=
	\norm{\int_{\RR^\dim} x (\d\al(x)-\d\be(x))  },
}
すなわちもはやノルムではない。これは、エネルギー距離ではWasserstein距離と同じ$f_{s\sharp}$下での線形スケーリングを得ることができないことを示している。
%
\todoK{compare the sqrt kernel associated to energy and sobolev and comment its slow decay}
\end{example}





%%%%%%%%%%%%%%%%%%%%%%%%%%%%%%%%%%%%%%%%%%%%%%%%%%%%%%%%%%%%%%%%%%%%%%%%%%%%%%%%%%%%%%%%%%%%%%%%%%%%
%%%%%%%%%%%%%%%%%%%%%%%%%%%%%%%%%%%%%%%%%%%%%%%%%%%%%%%%%%%%%%%%%%%%%%%%%%%%%%%%%%%%%%%%%%%%%%%%%%%%
%%%%%%%%%%%%%%%%%%%%%%%%%%%%%%%%%%%%%%%%%%%%%%%%%%%%%%%%%%%%%%%%%%%%%%%%%%%%%%%%%%%%%%%%%%%%%%%%%%%%
\section{Wasserstein空間はHilbert的ではない}\label{sec-non-embeddability}

\S\ref{sec:specialcases}で先に概説したWasserstein幾何の特殊な場合のいくつかは、最適輸送距離がときに閉じた形で計算できるという事実を明らかにした。またそのような場合に最適輸送距離が確率測度間の\emph{Hilbert的}計量であること、すなわち以下に定義するように、入力測度の空間からHilbert空間への写像$\tphi$が存在することも示した。

\begin{defn}\label{def-Hilbertmetric}
集合$\Z\times \Z$上で定義された距離$d$がHilbert的であるとは、Hilbert空間$\RKHS$と写像$\tphi:\Z\rightarrow \RKHS$が存在して、$\Z$の任意の対$z,z'$に対して$d(z,z')=\|\tphi(z)-\tphi(z')\|_\RKHS$が成り立つことをいう。
\end{defn}

例えば、注釈~\ref{rem-1d-ot-generic}は、Wasserstein計量が一変量分布間のHilbertノルムであることを示している。これは単に$\tphi$を測度にその一般化分位関数を対応させる写像として定義することによる。注釈~\ref{rem-dist-gaussians}は、本章の~\eqref{eq-univariate-gauss}に記されているように、一変量Gauss分布に対して二つの一変量Gauss分布間のWasserstein距離がそれらの平均と標準偏差の間のEuclid距離に過ぎないことを示している。

Hilbert的距離はデータ解析の文脈で用いる際に多くの好ましい性質を持つ~\citep{dattorro2010convex}。第一に、それらは容易に動径基底関数カーネルに変換できる。任意のHilbert的距離$\dist$に対して、$e^{-\dist^p/t}$は任意の$0\leq p\leq 2$と任意の正スカラー$t$に対して正定値カーネルであることが~\citep[Cor. 3.3.3, Prop. 3.2.7]{berg84harmonic}で示されている。Gaussカーネル（$p=2$）とラプラスカーネル（$p=1$）は、通常のEuclid距離を用いたこの結果の単純な応用である。カーネル法の全分野~\citep{Hofmann2008}は、正定値カーネル行列上で動作する凸学習アルゴリズムを定義するために、カーネル関数の正定値性に基づいている。Hilbert空間の点は、多次元尺度構成法~\citep{borg2005modern}のような単純な方法を用いて、低い歪み率で効率的に低次元に埋め込むこともできる~\citep{Johnson84}, \citep[\S V.6.2]{barvinok2002course}。

Hilbert的距離がこのような性質を持つため、Wasserstein距離が上で概説した以上のより一般的な設定、特に$\X$の次元が2以上の場合にもHilbert的であることが期待されるかもしれない。これは以下の同値性を用いて反証できる。

\begin{prop}\label{prop-negative-definite} 距離$d$がHilbert的であることと$d^2$が負定値であることは同値である。
\end{prop}
\begin{proof} 距離がHilbert的であれば、$d^2$が負定値であることは自明である。実際、$\Z$の$n$個の点が与えられると、和$\sum r_i r_j d^2(z_i,z_j)$は$\sum r_i r_j \|\tphi(z_i)-\tphi(z_j)\|_\RKHS^2$と書き直すことができ、$\sum r_i=0$であることを利用して展開すると$-2\sum r_i r_j \dotp{\tphi(z_i)}{\tphi(z_j)}_\RKHS$となり、Hilbert内積の定義により負となる。逆に$d^2$が負定値であれば、$d$がHilbert的であることは\citet{schoenberg38}による重要な結果（\cite[p. 82, Prop. 3.2]{berg84harmonic}に概説）から得られる。
\end{proof}

したがって、以下の命題で述べるように、Wasserstein距離の二乗が負定値でないことを示せば、それがHilbert的でないことを示すのに十分である。

\begin{prop}\label{prop-negative-definite-negative} $\X=\RR^\dim$で$\dim\geq 2$とし、基底コストを$\dist(x,y)=\norm{x-y}_2$とするとき、$p$-Wasserstein距離は$p=1,2$に対してHilbert的でない。
\end{prop}

\begin{proof} $\dim=2$に対して結果を証明すれば十分である。なぜなら、その次元での反例は任意のより高い次元での反例を得るのに十分だからである。$\RR^{2}$の4つのベクトル$x^1, x^2,x^3,x^4$上に台を持つ測度を用いた非確率的な反例を以下のように提示する。$x^{1}=[0,0], x^{2}=[1,0], x^{3}=[0,1], x^{4}=[1,1]$と定義する。次に、4次元の単体上の正則格子のすべての点を$1/4$刻みで考える。単体上にはそのような点が$35=\bigl(\!\bigl(\substack{4\\4}\bigl)\!\bigl)=\binom{4+4-1}{4}$個ある。格子上の各確率ベクトル$\a^i$は、$j\leq 4$に対して$\a^i_j$が集合$\{0,\tfrac{1}{4},\tfrac{1}{2},\tfrac{3}{4},1\}$に属し、$\sum_{j=1}^4\a^i_j=1$を満たす。与えられた$p$に対して、これらのヒストグラム間の$35\times 35$対Wasserstein距離行列$\D_p$が計算できる。$\D_p$が負定値でないことは、その要素ごとの二乗$\D_p^2$について$\J \D_p^2 \J$が正の固有値を持つことと同値であり、ここで$\J$は中心化行列$\J=\Identity_n-\tfrac{1}{n}\ones_{n,n}$である。これは図~\ref{fig-counterexample}に示されている通りである。
\end{proof}

\begin{figure}[h!]
\centering
\includegraphics[width=.6\linewidth]{counterexample_nsdness/maxeig_vs_p}
\caption{\label{fig-counterexample}
距離がHilbert的で\emph{ない}ことは、点集合の対二乗距離行列$\D_p^2$に対応する中心化行列$\J \D_p^2 \J$のスペクトルを見ることで示せる。距離がHilbert的であれば、そのような行列のスペクトルは必然的に非正である。ここでは命題~\ref{prop-negative-definite-negative}の証明で選ばれた点に対して、その行列の最大固有値の値を$p$の種々の値に対してプロットしている。得られた最大固有値を示す。これらの固有値はすべて正であり、これはこれらの$p$のすべての値に対して$p$-Wasserstein距離がHilbert的でないことを示している。
}
\end{figure}



%%%%%%%%%%%%%%%%%%%%%%%%%%%%%%%%%%%%%%%%%%%%%%%%%%%%%%%%%%%%%%%%%%%%%%%%%%%%%%%%%%%%%%%%%%%%%%%%%%%%
\subsection{埋め込みと歪み}

重要な研究の蓄積が、Hilbert的埋め込みを用いたWasserstein距離の近似の困難さを定量化している。
%
$\X=\RR^\dim$で$\dim\geq 2$の場合、測度を$\ell_2$空間に埋め込むと必然的に重要な\todoK{need to be more precise here}歪みが生じることが示されている~(\citet{naor-2005,andoni2015snowflake})。

$0 < p \leq 1$に対して$p$-Wasserstein空間を$\ell_1$に準等長的に埋め込むことは可能であるが（~\citep{indyk,andoni2008earth,do2011sublinear}を参照）、距離間の同値定数は次元$\dim$とともに急速に増大する。なお、$p=1$の場合、埋め込みは離散測度に対してのみ成り立つ（すなわち、埋め込み定数はスパイク間の最小距離に依存する）。\todoK{need to be more precise here}
%
密接に関連する埋め込み手法として、$\Wass_1$のLipschitz関数$f$の双対としての特徴付け（\S\ref{sec-w1-eucl}を参照）を用いて、Lipschitz制約$\norm{\nabla f}_1 \leq 1$をウェーブレット係数上の重み付き$\ell_1$球で近似する方法がある。~\citep{shirdhonkar2008approximate}を参照されたい。
%
このウェーブレット係数の重み付き$\ell_1$球は、いわゆる負の指数のベソフ空間を定義する~\citep{leeb2016holder}。
%
これらの埋め込み結果は、ダイアディック分割を用いたWasserstein距離の上界とも類似している。~\citep[Prop. 1]{weed2017sharp}および~\citep{fournier2015rate}も参照されたい。
%
これは$\ell_1$への準等長的な埋め込みも提供し（この埋め込みはスケーリングされたウェーブレット係数によって与えられる）、入力測度が一様格子上で離散化されている場合にこの埋め込みを線形時間で近似的に計算できるという利点がある。ウェーブレットの詳細は~\citep{mallat2008wavelet}を参照されたい。なお、Wasserstein的距離を計算するためにマルチスケール埋め込みを用いるアイデアは、コンピュータビジョンで広く用いられてきた。例えば~\citep{ling2006diffusion,grauman2005pyramid,CuturiNIPS2006,lazebnik2006beyond}を参照されたい。

%%%%%%%%%%%%%%%%%%%%%%%%%%%%%%%%%%%%%%%%%%%%%%%%%%%%%%%%%%%%%%%%%%%%%%%%%%%%%%%%%%%%%%%%%%%%%%%%%%%%
\subsection{最適輸送の負定値/正定値変種}

後に\S\ref{sec-sliced}で示すように、Wasserstein距離の\emph{スライス}近似、すなわち直線上に射影された測度のランダムな押し出しに対して計算される1次元方向輸送距離の和は、負定値関数の和として負定値である~\citep[\S 3.1.11]{berg84harmonic}。この結果は正定値カーネルを定義するために用いることができる~\citep{kolouri2016sliced}。正定値カーネルを回復する別の方法は、\citet{kosowsky1994invisible}がエントロピー正則化を導入するために提案したように、最適輸送問題を（可能なすべての輸送テーブル上の）最小問題ではなくソフトミン問題として定式化することである。そのソフトミンは、その負の指数関数（母関数とも呼ばれる）が正定値である項を定義する~\citep{cuturi2012positivity}。





%%%%%%%%%%%%%%%%%%%%%%%%%%%%%%%%%%%%%%%%%%%%%%%%%%%%%%%%%%%%%%%%%%%%%%%%%%%%%%%%%%%%%%%%%%%%%%%%%%%%
%%%%%%%%%%%%%%%%%%%%%%%%%%%%%%%%%%%%%%%%%%%%%%%%%%%%%%%%%%%%%%%%%%%%%%%%%%%%%%%%%%%%%%%%%%%%%%%%%%%%
%%%%%%%%%%%%%%%%%%%%%%%%%%%%%%%%%%%%%%%%%%%%%%%%%%%%%%%%%%%%%%%%%%%%%%%%%%%%%%%%%%%%%%%%%%%%%%%%%%%%
\section{OT、MMD、$\phi$-ダイバージェンスの経験的推定量}
\label{sec-empirical-wass}

応用的な設定において、二つの入力測度$(\al,\be) \in \Mm_+^1(\X)^2$が与えられたとき、重要な統計的問題は、$\al$からのサンプル$(x_i)_{i=1}^n$と$\be$からのサンプル$(y_j)_{j=1}^m$のみを用いて（通常は未知の）ダイバージェンス$D(\al,\be)$を近似することである。これらのサンプルはそれぞれの分布から独立同分布に従って抽出されたものと仮定する。

%%%%%%%%%%%%%%%%%%%%%%%%%%%%%%%%%%%%%%%%%%%%%%%%%%%%%%%%%%%%%%%%%%%%%%%%%%
\subsection{OTとMMDの経験的推定量}

Wasserstein距離$\Wass_p$（~\ref{eq-defn-wass-dist}を参照）とMMDノルム（\S\ref{sec-dual-norms}を参照）の両方に対して、分布間の未知の距離の直接的な推定量は、経験的測度間で直接計算することであり、理想的には後者から前者への収束速度を制御できることが望ましい。
\eq{
	D(\al,\be) \approx D(\hat \al_n,\hat \be_m) \qwhereq
	\choice{
		\hat \al_n \eqdef \frac{1}{n}\sum_i \de_{x_i},\\
		\hat \be_m \eqdef \frac{1}{m}\sum_j \de_{y_j}.
	}
}
ここで$\hat \al_n$と$\hat \be_m$はともに確率的測度であるため、$D(\hat \al_n,\hat \be_m)$は確率変数であることに注意されたい。
%
簡単のため、$\X$はコンパクトであると仮定する（非有界領域を扱うには入力測度のモーメントに追加の制約が必要となる）。

弱収束を距離化する（定義~\ref{dfn-weak-conv}を参照）このような双対距離に対して、弱収束$\hat\al_n \rightarrow \al$が成り立つので、$n \rightarrow +\infty$のとき$D(\hat \al_n,\hat \be_n) \rightarrow D(\al,\be)$となる。
%
しかし重要な問題は$D(\hat \al_n,\hat \be_n)$の$D(\al,\be)$への収束速度であり、この速度はしばしば$D$の「サンプル計算量」と呼ばれる。

$D(\al,\be) = \norm{\cdot}_{\TV}$の場合、TVノルムは弱収束を距離化しないため、$\norms{\hat \al_n - \hat \be_n}_{\TV}$は一致推定量ではない。すなわち、$\norm{\al-\be}_{\TV}$に収束しない。実際、確率1で$\norms{\hat \al_n - \hat \be_n}_{\TV}=2$となる。二つの離散測度の台が重ならないからである。同様の問題は他の$\phi$-ダイバージェンスでも生じ、経験的分布間のダイバージェンスを用いて推定することはできない。


%%%%
\paragraph{OTの収束速度}

$\X=\RR^\dim$で有界領域上に台を持つ測度に対して、$\dim>2$かつ$1 \leq p < +\infty$のとき、~\citep{dudley1969speed}により以下が示されている。
\eq{
	\EE( |\Wass_p(\hat \al_n,\hat \be_n)-\Wass_p(\al,\be)| ) = O(n^{-\frac{1}{\dim}}),
}
ここで期待値$\EE$は確率的サンプル$(x_i,y_i)_i$に関してとられている。この速度は、二つの測度の一方がLebesgue測度に関して密度を持つ場合、$\RR^\dim$では最適である。この結果は被覆数の概念を用いて一般の距離空間に対して証明され~\citep{dudley1969speed}、後に特に$\X=\RR^d$の場合に精緻化された~\citep{dereich2013constructive,fournier2015rate}。
% Proof on $\Wass_1$ implies on $\Wass_p$ for $1 \leq p < +\infty$ by monotonicity in $p$.
この速度は、測度が低次元部分領域上に台を持つ場合に精緻化される。~\citet{weed2017sharp}は、実際に速度が台の内在的次元に依存することを示した。\citeauthor{weed2017sharp}はまた、離散的に近似された測度（例えば小さい分散のGauss混合）のような収束の非漸近的挙動も研究した。
%
$\Wass_p(\hat \al_n,\hat \be_n)$のその平均$\Wass_p(\al,\be)$周りへの集中も証明することが可能である。~\citep{bolley2007quantitative,boissard2011simple,weed2017sharp}を参照されたい。

%%%
\paragraph{MMDの収束速度}

~\eqref{eq-kernel-dual}で定義されたRKHSノルムの双対である弱ノルム$\norm{\cdot}^2_{\Krkhs}$（MMDとも呼ばれる）に対して、Wasserstein距離とは対照的に、サンプル計算量は周囲の次元に依存しない。
\eq{
	\EE( |\norms{\hat \al_n-\hat \be_n}_{\Krkhs} - \norm{\al-\be}_{\Krkhs} | ) = O(n^{-\frac{1}{2}});
}
~\citep{sriperumbudur2012empirical}を参照されたい。
%
図~\ref{fig-sample-complexity}はWasserstein距離とMMD距離のサンプル計算量の速度の数値的比較を示している。
%
ただし、$\norms{\hat \al_n-\hat \be_n}_{\Krkhs}^2$は$\norm{\al-\be}_{\Krkhs}^2$のわずかに偏った推定量であることに注意されたい。不偏推定量を定義し、例えばこのような損失の最小化にSGDを使用できるようにするためには、以下の不偏推定量を用いるべきである。
\begin{align*}
	\text{MMD}_{\Krkhs}(\hat \al_n,\hat \be_n)^2 & \eqdef
		\frac{1}{n(n-1)} \sum_{i,i'} \Krkhs(x_i,x_{i'})	+
		\frac{1}{n(n-1)}\sum_{j,j'} \Krkhs(y_j,y_{j'}) \\
		 & \qquad - 2
		\frac{1}{n^2}\sum_{i,j} \Krkhs(x_i,y_j),
\end{align*}
これは~\eqref{eq-mmd-discr}と比較されるべきである。これは$\EE(\text{MMD}_{\Krkhs}(\hat \al_n,\hat \be_n)^2) = \norm{\al-\be}_{\Krkhs}^2$を満たす。~\citep{gretton2012kernel}を参照されたい。


\begin{figure}[h!]
\centering
\begin{tabular}{@{}c@{\hspace{1mm}}c@{}}
\includegraphics[width=.49\linewidth]{sample-complexity/energy-rates-cropped}&
\includegraphics[width=.49\linewidth]{sample-complexity/wass2-rates-cropped}\\
エネルギー距離 $\norm{\cdot}_{H^{-1}}$ & $\Wass_2$
\end{tabular}
\caption{\label{fig-sample-complexity}
$D$が例~\ref{exp-energy-dist}で定義されたエネルギー距離$D=\norm{\cdot}_{H^{-1}}$（すなわち$H^{-1}$ノルム）（左）およびWasserstein距離$D=\Wass_2$（右）の場合の$\log_{10}(n)$の関数としての$\log_{10}(D(\hat\al_n,\hat\al_n'))$の減衰。
%
ここで$(\hat\al_n,\hat\al_n')$は$\al$の二つの独立な経験的分布であり、$\al$は単位立方体$[0,1]^\dim$上の一様分布であり、$\dim \in \{2,3,5\}$のいくつかの値に対して検証されている。
%
網掛け部分は$\log(D(\hat\al_n,\al))$の標準偏差$\pm$での信頼区間を示している。
}
\end{figure}



%%%%%%%%%%%%%%%%%%%%%%%%%%%%%%%%%%%%%%%%%%%%%%%%%%%%%%%%%%%%%%%%%%%%%%%%%%
\subsection{$\phi$-ダイバージェンスの経験的推定量}

離散サンプルから~\eqref{def_divergence}で定義された$\Divergm_\phi(\al|\be)$を$\Divergm_\phi(\hat \al_n|\hat \be_n)$を用いて近似することはできない。実際、この量は$+\infty$であるか（例えば$\KL$ダイバージェンスの場合）、あるいは$n \rightarrow +\infty$のとき$\Divergm_\phi(\al|\be)$に収束しない（例えば$\TV$ノルムの場合）。
%
代わりに、離散的な経験的測度を何らかの方法で平滑化して密度に置き換えるために、密度推定量を用いる必要がある。~\citep{silverman1986density}を参照されたい。
%
Euclid空間$\Xx=\RR^\dim$において、滑らかな窓関数と帯域幅$\si>0$を用いて$h_\si=h(\cdot/\si)$を導入すると、$\al$に対する密度推定量はこのカーネルとの畳み込みによって定義される。
\eql{\label{eq-kernel-density-estimate}
	\hat \mu_n \star h_\si = \frac{1}{n}\sum_i h_\si(\cdot-x_i).
}
%
すると$\phi$ダイバージェンスを以下で近似できる。
\eq{
	\Divergm_\phi^\si(\hat \al_n |\hat \be_n)
	\eqdef
	\frac{1}{n} \sum_{j=1}^n
	\phi\pa{
		\frac{
			\sum_i h_\si(y_j-x_i)
		}{
			\sum_{j'} h_\si(y_j-y_{j'}),
		}
	}
}
ここで$\si$はサンプル数$n$と次元$\dim$に適合させる必要がある。
%
固定帯域幅$\si$の選択を回避し、代わりに最近傍数$k$を選択するノンパラメトリック推定量を考案することも可能である。これらの手法は典型的に最近傍間の距離を利用し~\citep{loftsgaarden1965nonparametric}、これは帯域幅$\si$を局所サンプリング密度に局所的に適応させることに類似する。$\De_k(x)$を$x \in \RR^\dim$と$(x_i)_{i=1}^n$の中のその第$k$最近傍との距離とすると、密度推定量は以下で定義される。
\eql{\label{eq-knn-density-estimate}
	\rho_{\hat \al_n}^k(x) \eqdef \frac{k/n}{ |B_\dim| \De_k(x)^r },
}
ここで$|B_\dim|$は$\RR^\dim$における単位球の体積である。
%
~\eqref{eq-kernel-density-estimate}において幅$\si$の領域内に落ちるサンプル数を何らかの方法で「数える」代わりに、この式~\eqref{eq-knn-density-estimate}は$k$個のサンプルを包含するのに必要な半径を推定する。図~\ref{fig-density-estimation}は推定量~\eqref{eq-kernel-density-estimate}と~\eqref{eq-knn-density-estimate}を比較している。
%
典型的な応用例は~\eqref{eq-discr-entropy}でエントロピー汎関数（Lebesgue測度に関するKLダイバージェンス）について詳述されている。詳細は~\citep{BerishaHero}を参照されたい。


\newcommand{\myFigDE}[1]{\includegraphics[width=.32\linewidth]{density-estimation/#1}}

\begin{figure}[h!]
\centering
\begin{tabular}{@{}c@{\hspace{1mm}}c@{\hspace{1mm}}c@{}}
\myFigDE{kernel-1} & \myFigDE{kernel-3} & \myFigDE{kernel-5} \\
$\si=2.5 \cdot 10^{-3}$ & $\si=15 \cdot 10^{-3}$ & $\si=25 \cdot 10^{-3}$ \\
\myFigDE{nn-1} & \myFigDE{nn-3} & \myFigDE{nn-5} \\
$k=1$ & $k=50$ & $k=100$
\end{tabular}
\caption{\label{fig-density-estimation}
二つのGauss分布の混合からの$n=200$サンプルに対する、カーネル密度推定$\hat \mu_n \star h_\si$（上段、Gaussカーネル$h$を使用）と$k$-最近傍推定$\rho_{\hat \al_n}^k$（下段）の比較。
}
\end{figure}



%%%%%%%%%%%%%%%%%%%%%%%%%%%%%%%%%%%%%%%%%%%%%%%%%%%%%%%%%%%%%%%%%%%%%%%%%%%%%%%%
%%%%%%%%%%%%%%%%%%%%%%%%%%%%%%%%%%%%%%%%%%%%%%%%%%%%%%%%%%%%%%%%%%%%%%%%%%%%%%%%
%%%%%%%%%%%%%%%%%%%%%%%%%%%%%%%%%%%%%%%%%%%%%%%%%%%%%%%%%%%%%%%%%%%%%%%%%%%%%%%%
\section{エントロピー正則化：OTとMMDの間}
\label{sec-entropy-ot-mmd}

命題~\ref{prop-sinkhorn-div}に従い、Sinkhornダイバージェンスは以下で定義されることを想起する。
\eq{
	\SINKHORNP_\C^\varepsilon(\a,\b) \eqdef
	\dotp{\P^\star}{\C} =
	\dotp{e^{\frac{\fD^\star}{\varepsilon}}}{(\K\odot \C)e^{\frac{\gD^\star}{\varepsilon}}},
}
ここで$\P^\star$は~\eqref{eq-regularized-discr}の解であり、$(\fD^\star,\gD^\star)$は~\eqref{eq-dual-formulation}の解である。
%
ある距離$\dist$に対して$\C_{i,j}=\dist(x_i,x_j)^p$と仮定すると、~\eqref{eq-pair-discr}の形の二つの離散確率分布に対して、これは正則化されたWassersteinコストを定義する。
\eq{
	\Wass_{p,\varepsilon}(\al,\be)^p \eqdef \SINKHORNP_\C^\varepsilon(\a,\b).
}
この定義は任意の入力分布（必ずしも離散的でなくてよい）に以下のように一般化される。
\eq{
	\Wass_{p,\varepsilon}(\al,\be)^p \eqdef \int_{\X \times \X} \dist(x,y)^p \d\pi^\star(x,y),
}
ここで$\pi^\star$は~\eqref{eq-entropic-generic}の解である。

正則化によって導入されるバイアス（特に$\Wass_{p,\varepsilon}(\al,\al) \neq 0$）を打ち消すために、補正された正則化ダイバージェンスを導入する。
\eq{
	\tilde\Wass_{p,\varepsilon}(\al,\be)^p \eqdef
	2 \Wass_{p,\varepsilon}(\al,\be)^p
	-
	\Wass_{p,\varepsilon}(\al,\al)^p
	-
	\Wass_{p,\varepsilon}(\be,\be)^p.
}
$e^{-c/\epsilon}$が正定値カーネルであれば、関連する補正されたダイバージェンス（$\SINKHORNP_\C^\varepsilon$の代わりに$\MKD_\C^\varepsilon$を用いて得られるもの）は正であることが~\citep{feydy2018interpolating}で証明されている。
%
例えば~\citet{amari2017information}が提案しているように、正則化最適輸送を用いた他の再正規化スキームを定義することも可能であることに注意されたい。

\begin{figure}[h!]
\centering
\begin{tabular}{@{}c@{\hspace{1mm}}c@{}}
\includegraphics[width=.49\linewidth]{sample-complexity-regul/influ-epsilon-d2-cropped}&
\includegraphics[width=.49\linewidth]{sample-complexity-regul/influ-epsilon-d5-cropped}\\
$\dim=2$ & $\dim=5$
\end{tabular}
\caption{\label{fig-sample-complexity-regul}
種々の$\varepsilon$に対する$p=3/2$での$\EE(\log_{10}(\tilde\Wass_{p,\varepsilon}(\hat\al_n,\hat\al_n')))$の$\log_{10}(n)$の関数としての減衰。ここで$\al$は図~\ref{fig-sample-complexity}と同じものである。\todoK{redo the plot with higher precision}
}
\end{figure}


以下の命題は~\citep{ramdas2017wasserstein}に証明があり、この正則化ダイバージェンスが例~\ref{exp-energy-dist}で定義されたWasserstein距離とエネルギー距離の間を補間することを示している。

\begin{prop}
	以下が成り立つ。
	\eq{
		\tilde\Wass_{p,\varepsilon}(\al,\be)  \overset{\varepsilon \rightarrow 0}{\longrightarrow} 2 \Wass_{p}(\al,\be)
		\qandq
		\tilde\Wass_{p,\varepsilon}(\al,\be)^p \overset{\varepsilon \rightarrow +\infty}{\longrightarrow} \norm{\al-\be}_{\text{ED}(\X,\dist)}^2,
	}
	ここで$\norm{\cdot}_{\text{ED}(\X,\dist)}$は~\eqref{eq-defn-ed}で定義されている。
\end{prop}

図~\ref{fig-sample-complexity-regul}は$\epsilon$のサンプル計算量の速度への影響を数値的に示している。$\Xx=\RR^d$上で$c(x,y)=\norm{x-y}^2$の場合、これらの速度がOTとMMDの速度の間を補間することが~\cite{genevay2018sample}で証明されている。




